% Created 2025-01-17 vie 13:13
% Intended LaTeX compiler: pdflatex
\documentclass[11pt]{article}
\usepackage[utf8]{inputenc}
\usepackage[T1]{fontenc}
\usepackage{graphicx}
\usepackage{longtable}
\usepackage{wrapfig}
\usepackage{rotating}
\usepackage[normalem]{ulem}
\usepackage{amsmath}
\usepackage{amssymb}
\usepackage{capt-of}
\usepackage{hyperref}
\usepackage{minted}

\usepackage{parskip}
\usepackage[utf8]{inputenc} %% For unicode chars
\usepackage{placeins}
\usepackage[margin=2.50cm]{geometry}
\usepackage[T1]{fontenc}
\usepackage{mathpazo}
\linespread{1.05}
\usepackage[scaled]{helvet}
\usepackage{courier}
\hypersetup{colorlinks=true,linkcolor=blue}
\RequirePackage{fancyvrb}
\DefineVerbatimEnvironment{verbatim}{Verbatim}{fontsize=\small,formatcom = {\color[rgb]{0.5,0,0}}}
\usepackage{mdframed}
\BeforeBeginEnvironment{minted}{\begin{mdframed}}
\AfterEndEnvironment{minted}{\end{mdframed}}
\setcounter{secnumdepth}{3}
\author{José Manuel Sánchez Álvarez}
\date{}
\title{PHP-basics}
\hypersetup{
 pdfauthor={José Manuel Sánchez Álvarez},
 pdftitle={PHP-basics},
 pdfkeywords={},
 pdfsubject={},
 pdfcreator={Emacs 27.1 (Org mode 9.6.18)}, 
 pdflang={Spanish}}
\begin{document}

\maketitle
\setcounter{tocdepth}{3}
\tableofcontents

\setlength\parindent{10pt}

\section{PHP Basics}
\label{sec:org686c6b6}
Welcome to the first parte of my book, where I will teach you the basics
of PHP, a popular and powerful scripting language for web development.
PHP stands for PHP: Hypertext Preprocessor, which is a recursive acronym
that reflects how PHP can process text and embed code within HTML
documents. In this chapter, you will learn how to set up your
environment, write your first PHP script, and use variables and data
types in PHP.

\subsection{Setting up the environment}
\label{sec:orgf77d4e1}
Before you can start writing PHP code, you need to have a web server
that can run PHP scripts. There are many options for web servers, such
as Apache, Nginx, or IIS, but for simplicity, we will use a built-in web
server that comes with PHP. To use the built-in web server, you need to
have PHP installed on your computer. You can download PHP from the
official website (\url{https://www.php.net/downloads.php}) and follow the
instructions for your operating system. Once you have installed PHP, you
can open a terminal or command prompt and navigate to the folder where
you want to store your PHP files. Then, you can start the web server by
typing the following command:

\begin{minted}[]{bash}
php -S localhost:8000
\end{minted}

This will start a web server on your local machine, listening on port
\begin{enumerate}
\item You can access the web server by opening a browser and typing
\end{enumerate}
\url{http://localhost:8000} in the address bar. You should see a message
saying "No input file specified.", which means that the web server is
working but there is no PHP file to execute.

\subsection{Your first PHP script}
\label{sec:orgf14ac4c}
To create your first PHP script, you need to create a file with the .php
extension in the folder where you started the web server. For example,
you can create a file called index.php and open it with a text editor.
Then, you can write the following code in the file:

\begin{minted}[]{php}
<?php
echo "Hello, world!";
?>
\end{minted}

The \texttt{<?php} and \texttt{?>} tags indicate the start and end of a PHP code
block. The \texttt{echo} statement outputs a string to the browser. To run this
script, you need to save the file and refresh the browser. You should
see "Hello, world!"  displayed on the web page.

\subsection{Variables and Data Types}
\label{sec:orgce6c1df}
A variable is a container that holds a value that can change during the
execution of a script. In PHP, you can declare a variable by using the \$
sign followed by a name. For example:

\begin{minted}[]{php}
<?=
  $name = "Alice";
\end{minted}

This assigns the string "Alice" to the variable \$name. You can use
variables in expressions or output them using echo. For example:

\begin{minted}[]{php}
<?=
  echo "Hello, $name!";
\end{minted}

This will output "Hello, Alice!" to the browser.

PHP supports different types of data that can be stored in variables,
such as integers (whole numbers), floats (decimal numbers), strings
(text), booleans (true or false), arrays (collections of values),
objects (instances of classes), resources (handles to external
resources), and null (no value). You can check the type of a variable by
using the gettype() function. For example:

\begin{minted}[]{php}
<?=
  $age = 25;
  echo gettype($age);
\end{minted}

This will output "integer" to the browser.

PHP is a \emph{loosely typed} language, which means that you do not have to
specify the type of a variable when you declare it. PHP will
automatically convert the type of a variable depending on how you use
it. For example:

\begin{minted}[]{php}
<?=
  $age = 25;
  $age = "twenty-five";
  echo gettype($age);
\end{minted}

This will output "string" to the browser, because \$age was assigned a
string value.

\subsection{Variable scope}
\label{sec:orgb09a4d7}
The scope of a variable is the part of the script where the variable is
accessible. In PHP, there are three main scopes: global, local, and
static.

\begin{enumerate}
\item \textbf{Global scope}: A variable declared outside any function or class is in
the global scope and can be accessed from anywhere in the script.
However, to access a global variable from inside a function or class,
you need to use the global keyword or the \$GLOBALS array. For
example:

\begin{minted}[]{php}
<?=
  $x = 10; // global variable

  function test() {
    global $x; // access global variable using global keyword
    echo $x; // output 10
  }

  test();

  function test2() {
    echo $GLOBALS['x']; // access global variable using $GLOBALS array
  }

  test2();
\end{minted}

\item Local scope: A variable declared inside a function or class is in the
local scope and can only be accessed from within that function or
class. For example:

\begin{minted}[]{php}
<?=
  function test3() {
    $y = 20; // local variable
    echo $y; // output 20
  }

  test3();

  echo $y; // error: undefined variable
\end{minted}

\item Static scope: A static variable is declared inside a function using
the static keyword and retains its value between function calls. For
example:

\begin{minted}[]{php}
<?=
  function test4() {
    static $z = 0; // static variable
    $z++;
    echo $z; // output 1, 2, 3, ... on each function call
  }

  test4();
  test4();
  test4();
\end{minted}
\end{enumerate}

\subsection{Expressions, operators, operands}
\label{sec:org026340e}
An expression is a combination of values, variables, operators, and
functions that evaluates to a single value. For example:

\begin{minted}[]{php}
<?=
  $sum = 2 + 3; // expression that evaluates to 5
\end{minted}

An operator is a symbol that performs a specific operation on one or
more operands. An operand is a value or variable that is used by an
operator. For example:

\begin{minted}[]{php}
<?=
  2 + 3 // + is an operator, 2 and 3 are operands
\end{minted}

PHP supports various types of operators, such as arithmetic,
assignment, comparison, logical, string, array, and conditional
operators. You can find a complete list of operators and their
precedence and associativity rules in the \href{https://www.php.net/manual/en/language.operators.php}{PHP manual}

Some examples of operators and their usage are:

\begin{enumerate}
\item Arithmetic operators: Perform basic mathematical operations such as
addition (\texttt{+}), subtraction (\texttt{-}), multiplication (\texttt{*}), division
(\texttt{/}), modulus (\texttt{\%}), and exponentiation (\texttt{**}). For example:

\begin{minted}[]{php}
<?=
  $a = 10;
  $b = 3;
  echo $a + $b; // output 13
  echo $a - $b; // output 7
  echo $a * $b; // output 30
  echo $a / $b; // output 3.3333333333333
  echo $a % $b; // output 1
  echo $a ** $b; // output 1000
\end{minted}

\item Assignment operators: Assign a value to a variable or modify the
value of a variable by another value. The basic assignment operator
is (\texttt{=}), which assigns the right operand to the left operand. For
example:

\begin{minted}[]{php}
<?=
  $c = 5; // assign 5 to $c
\end{minted}

\item There are also compound assignment operators that combine an
arithmetic or string operator with the assignment operator. For
example:

\begin{minted}[]{php}
<?=
  $c += 2; // equivalent to $c = $c + 2;
  $c -= 2; // equivalent to $c = $c - 2;
  $c *= 2; // equivalent to $c = $c * 2;
  $c /= 2; // equivalent to $c = $c / 2;
  $c %= 2; // equivalent to $c = $c % 2;
  $c **= 2; // equivalent to $c = $c ** 2;
  $c .= "hello"; // equivalent to $c = $c . "hello";
\end{minted}

\item Comparison operators: Compare two values and return a boolean value
(\texttt{true} or \texttt{false}) based on the result of the comparison. There are
two types of comparison operators: loose and strict. Loose comparison
operators compare the values without checking their types, while
strict comparison operators compare both the values and their types.
For example:

\begin{minted}[]{php}
<?=
  $d = "10";
  $e = 10;

  echo ($d == $e); // output true (loose equality)
  echo ($d === $e); // output false (strict equality)
  echo ($d != $e); // output false (loose inequality)
  echo ($d !== $e); // output true (strict inequality)
  echo ($d < $e); // output false (less than)
  echo ($d > $e); // output false (greater than)
  echo ($d <= $e); // output true (less than or equal to)
  echo ($d >= $e); // output true (greater than or equal to)
\end{minted}

\item Logical operators: Perform logical operations on one or more boolean
values and return a boolean value based on the result of the
operation. There are four logical operators:
\texttt{and (\&\&), or (||), xor (\textasciicircum{})}, and \texttt{not (!)}. For example:

\begin{minted}[]{php}
<?=
  $f = true;
  $g = false;

  echo ($f && $g); // output false (logical and)
  echo ($f || $g); // output true (logical or)
  echo ($f ^ $g); // output true (logical xor)
  echo (!$f); // output false (logical not)
\end{minted}

\item String operators: Concatenate two strings or append one string to
another. There are two string operators: concatenation (.) and
concatenation assignment (.=). For example:

\begin{minted}[]{php}
<?=
  $h = "Hello";
  $i = "world";

  echo ($h . " " . $i); // output Hello world (concatenation)
  $h .= " "; // append a space to $h
  $h .= $i; // append $i to $h
  echo ($h); // output Hello world (concatenation assignment)
\end{minted}

\item Array operators: Combine two arrays or compare two arrays. There are
five array operators: union (\texttt{+}), equality (\texttt{==}), identity (\texttt{===}),
inequality (\texttt{!=}), and non-identity (\texttt{!==}). For example:

\begin{minted}[]{php}
<?=
  $j = array("a" => "apple", "b" => "banana");
  $k = array("c" => "cherry", "d" => "durian");

  print_r($j + $k); // output Array ([a] => apple [b] => banana
                    // [c] => cherry [d] => durian) (array union)
  echo ($j == $k); // output false (array equality)
  echo ($j === $k); // output false (array identity)
  echo ($j != $k); // output true (array inequality)
  echo ($j !== $k); // output true (array non-identity)
\end{minted}

\item Conditional operator: Also known as the ternary operator, it
evaluates a condition and returns one of two values based on whether
the condition is true or false. The syntax is \texttt{condition ?
   value\_if\_true : value\_if\_false}. For example:

\begin{minted}[]{php}
<?=
 $l = 10;
 $m = ($l > 5) ? "greater than 5" : "not greater than 5";
 echo $m; // Output: "greater than 5"
\end{minted}
\end{enumerate}
\end{document}
